\scanpage{18}
% archive foliation: 17
$\unsure{\xi|\zeta}$-fibré en droites projectives et $\varphi:P^1(\mathbb C)\longrightarrow P_x$ un iso. L'opération de $H$ dessus est \unsure{l'évidente}.

Or le groupe $\operatorname{Aut}(\pi)$ opère \inserted{à droite} sur $\uline{\operatorname{Hom}}(\pi,H)$ ; \deleted{et} si $u\in\operatorname{Aut}(\pi)$, $\rho\in\uline{\operatorname{Hom}}(\pi,H)$, on pose $\rho.u=\rho\circ u$ \inserted{(les op. commutent à celles de $H$)}. Si $u=\operatorname{int}(g)$ ($g\in\pi$), on a $\rho\circ u=\operatorname{int}(\rho(g))\circ\unsure{u|\rho}$, donc par passage au quotient : ${}_H\!\backslash\uline{\operatorname{Hom}}^!(\pi,H)\deleted{\unsure{/H}}$ c'est
\[
\unsure{\gamma|\Gamma}:\unsure{\gamma(\pi)|\Gamma(\pi)}=\operatorname{Aut\,ext}(\pi)=\operatorname{Aut}(\pi)/\underbrace{\operatorname{Aut\,int}(\pi)}_{\simeq\pi}
\]
qui opère sur lui. Par transport de structure, ce groupe opère sur $M_g^{\unsure{\xi|\zeta}}(X)$, mais sans doute cette opération n'est pas algébrique.
\margin{$?\ \big|$}

Si $Y$ est une courbe analytique compacte \inserted{connexe} de genre $g$, la donnée d'un iso. ext.
\[
\pi_1(Y)\simeq\pi=\pi_1(X,x)
\]
définit, via la représentation \uline{canonique}
\[
\rho_Y:\pi_1(Y)\longrightarrow H
\]
(définie mod autom. int. dans $H$), un élément de $M_g^{\unsure{\xi|\zeta}}(X)$. On trouve ainsi, si $M_g(\pi)\inserted{=M_g(X)}$ est l'espace de Teichmüller (à rigidifications définies par $\pi$), on trouve un morphisme analytique \uline{injectif}
\diagram{\begin{tikzcd}
M_g(X)\arrow[r,hook]&M_g^{\unsure{\xi|\zeta}}(X)
\end{tikzcd}}
qui commute aux \unsure{op.} du gr. de Teichmüller $\unsure{\gamma|\Gamma}$. On suspecte que l'image \deleted{n'est}\inserted{\unsure{ne soit|ne serait}} pas un sous-ens. constructible du pt de vue algébrique. Elle est formée des hom \uline{injectifs} $\pi\lhook\joinrel\xrightarrow{\rho}H$ pour lesquels il existe un domaine homogène $D\subset P^1(\mathbb C)$ inv. par $\rho$ (il sera unique) et que \deleted{$D/\unsure{Y}$ soit \uline{compact}} \inserted{$\pi$ y} opère librement et avec un quotient compact.

On définit aussi $M_{\mathrm{conn}}(X)$, qui classifie les \deleted{connexions projectives}\inserted{\unsure{$P(1)$-fibrés stratifiés}} sur $X$, \uline{munis} d'une section nulle, \unsure{non stationnaire}. \deleted{Il est} C'est un torseur sur le vectoriel de dim finie $H^0(X,\unsure{\underline{\omega}|\omega}^{\otimes2}_{X/\mathbb C})$, et on a
\diagram{\begin{tikzcd}
M_{\uline{\mathrm{conn}}}(X)\arrow[r]&M_g^{\unsure{\xi|\zeta}}(X)
\end{tikzcd}}
\margin{Attention, le $\unsure{\xi|\zeta}$-fibré défini par une \unsure{conn. proj.} \unsure{peut avoir des autom.}}
\deleted{\unsure{Ceci}} est un morphisme \uline{an.} (\deleted{\unsure{et localt injectif}})
(\uline{NB} C'est en fait un morphisme de schémas algébriques sur $\mathbb C$)
\deleted{\unsure{Et il est}}\inserted{En fait c'est un} \uline{mono} -- i.e. deux connexions projectives \unsure{différentes|définies}, sur $X$ \inserted{ne} \unsure{peuvent}\deleted{\unsure{alors}} donner des $\unsure{\xi|\zeta}$-fibrés isomorphes.
\margin{$\forall$ $\unsure{\xi|\zeta}$-iso \unsure{respecte la section marquée}.}

\inserted{En effet} le fibré \deleted{\unsure{On a d'ailleurs déjà} \illegible{1}} associé aux deux conn. proj. est le même $P=P(P^1_{X/S}(\unsure{\underline{\omega}|\omega}))$ et un $\unsure{\xi|\zeta}$-iso. est (\deleted{\unsure{donc}} un automorphisme de $P(P^1_{X/S}(\unsure{\underline{\omega}|\omega}))$ donc défini par un $L$ inv. et un iso.
\[
P^1_{X/S}(\unsure{\underline{\omega}|\omega})\ (\unsure{\text{noté }E})\simeq P^1_{X/S}(\unsure{\underline{\omega}|\omega})\otimes L
\]
i.e. $E\simeq E\otimes L$. On aura donc $\det E\simeq\det E\otimes L^2$ i.e. $L^2\simeq\mathcal O_X$, et un épi $E\longrightarrow L$ i.e. $P^1_{X/S}(\unsure{\underline{\omega}|\omega})\longrightarrow L$. Or, \deleted{\unsure{d'un tel}} \deleted{\unsure{par déf. nul par un}} $\unsure{\underline{\omega}|\omega}\longrightarrow L$ \deleted{\unsure{On a}} morphisme induit est une section de $\unsure{\underline{\omega}|\omega}^{-1}\otimes L$, or $H^0(\unsure{\underline{\omega}|\omega}^{-1}\otimes L)=0$ car
